% GENERATED FILE. Edit canonical lessons, then run npm run generate:books.
% canonical-source-hash: fnv1a64:dec0866a7f6eae97
% canonical-lessons: AR-C303-ustura, AR-C303-khurafa, AR-C303-batal, AR-C303-shirrir, AR-C303-takhayyala

\chapter{Legends and heroes}
\label{ch:ar-legends-and-heroes}
\hlchaptermodality{303}

\begin{chapteropening}
\textbf{By the end of this chapter:} \emph{I can say a legend, a myth, a hero, a villain and to imagine.}

Complete the last lesson of chapter 303: I can say a legend, a myth, a hero, a villain and to imagine.
\end{chapteropening}

\section[\texorpdfstring{\ar{أسطورة} (usṭūra)}{أسطورة (usṭūra)}]{\ar{أسطورة} (usṭūra) — a legend}

\label{lesson:AR-C303-ustura}



\begin{quote}
Before the new one: say the Arabic for to host, then the Arabic for a spokesperson.
\end{quote}

\subsection*{You'll want to know: \ar{أسطورة}}
\textbf{\ar{أسطورة}} — \emph{usṭūra} — \textquotedblleft{}a legend\textquotedblright{}.

\begin{quote}
\ar{سمعنا} \ar{أسطورة} \ar{عن} \ar{جبل} \ar{قديم}.
\end{quote}

\begin{quote}
\emph{samiʿnā usṭūra ʿan jabal qadīm.}
\end{quote}

\begin{quote}
We heard a legend about an old mountain.
\end{quote}

Plural: \ar{أساطير}, \emph{asāṭīr}.

\begin{cousinweb}
\textbf{\ar{أسطورة}} (\emph{usṭūra}) is filed by Lane under the root \emph{s-ṭ-r}, \textquotedblleft{}a line of writing\textquotedblright{}, and its plural \emph{asāṭīr} is old. It was borrowed from Greek \emph{historia}, \textquotedblleft{}an inquiry, a story\textquotedblright{} (EtymDB); others see a native word, so the origin is not settled.
\end{cousinweb}

\begin{grammarlens}[title={What you've built}]
One more word for telling what happened.
\end{grammarlens}

\subsection*{Your turn}
\begin{itemize}
\raggedright
  \item \emph{Say it:} \emph{usṭūra}
  \item \emph{Say it:} \emph{usṭūra}, once more
  \item \emph{Say it:} \emph{samiʿnā usṭūra ʿan jabal qadīm.} — We heard a legend about an old mountain.
  \item \emph{Recall:} say the Arabic for to host, then the Arabic for a spokesperson, then say \emph{usṭūra} again
\end{itemize}

\begin{tcolorbox}[breakable,colback=teal!4,colframe=teal!35!black,title={Before you move on}]
What does \ar{أسطورة} mean? (\textquotedblleft{}A legend\textquotedblright{}.) Say it once more.
\end{tcolorbox}
\section[\texorpdfstring{\ar{خرافة} (khurāfa)}{خرافة (khurāfa)}]{\ar{خرافة} (khurāfa) — a myth, a superstition}

\label{lesson:AR-C303-khurafa}



\begin{quote}
Before the new one: say the Arabic for a spokesperson, then the Arabic for a legend.
\end{quote}

\subsection*{You'll want to know: \ar{خرافة}}
\textbf{\ar{خرافة}} — \emph{khurāfa} — \textquotedblleft{}a myth, a superstition\textquotedblright{}.

\begin{quote}
\ar{قال} \ar{جدي} \ar{إن} \ar{القصة} \ar{خرافة}.
\end{quote}

\begin{quote}
\emph{qāla jaddī inna al-qiṣṣa khurāfa.}
\end{quote}

\begin{quote}
My grandfather said the story was a myth.
\end{quote}

Plural \ar{خرافات}. Of a belief it means a superstition.

\begin{cousinweb}
\textbf{\ar{خرافة}} (\emph{khurāfa}) comes from a name, as Lane records under the root \emph{kh-r-f} (Lane): Khurāfa, a man of the tribe of ʿUdhra, told what he had seen among the jinn and was called a liar, and \emph{khurāfāt} became tales deemed pretty but untrue.
\end{cousinweb}

\begin{grammarlens}[title={What you've built}]
One more word for telling what happened.
\end{grammarlens}

\subsection*{Your turn}
\begin{itemize}
\raggedright
  \item \emph{Say it:} \emph{khurāfa}
  \item \emph{Say it:} \emph{khurāfa}, once more
  \item \emph{Say it:} \emph{qāla jaddī inna al-qiṣṣa khurāfa.} — My grandfather said the story was a myth.
  \item \emph{Recall:} say the Arabic for a spokesperson, then the Arabic for a legend, then say \emph{khurāfa} again
\end{itemize}

\begin{tcolorbox}[breakable,colback=teal!4,colframe=teal!35!black,title={Before you move on}]
What does \ar{خرافة} mean? (\textquotedblleft{}A myth, a superstition\textquotedblright{}.) Say it once more.
\end{tcolorbox}
\section[\texorpdfstring{\ar{بطل} (baṭal)}{بطل (baṭal)}]{\ar{بطل} (baṭal) — a hero; a champion}

\label{lesson:AR-C303-batal}



\begin{quote}
Before the new one: say the Arabic for a legend, then the Arabic for a myth.
\end{quote}

\subsection*{You'll want to know: \ar{بطل}}
\textbf{\ar{بطل}} — \emph{baṭal} — \textquotedblleft{}a hero; a champion\textquotedblright{}.

\begin{quote}
\ar{كان} \ar{جدي} \ar{بطلا} \ar{في} \ar{قريته}.
\end{quote}

\begin{quote}
\emph{kāna jaddī baṭalan fī qaryatihi.}
\end{quote}

\begin{quote}
My grandfather was a hero in his village.
\end{quote}

Plural: \ar{أبطال}, \emph{abṭāl}.

In sport \ar{بطل} is a champion, and the hero of a film or novel is its \ar{بطل} too.

\begin{cousinweb}
\textbf{\ar{بطل}} (\emph{baṭal}) is built on the native Arabic root \emph{b-ṭ-l} (Lane), \textquotedblleft{}to come to nothing\textquotedblright{}. Lane explains \emph{baṭal}, a brave man, as one whose wound goes for nothing, so that it does not hold him back.
\end{cousinweb}

\begin{grammarlens}[title={What you've built}]
One more word for telling what happened.
\end{grammarlens}

\subsection*{Your turn}
\begin{itemize}
\raggedright
  \item \emph{Say it:} \emph{baṭal}
  \item \emph{Say it:} \emph{baṭal}, once more
  \item \emph{Say it:} \emph{kāna jaddī baṭalan fī qaryatihi.} — My grandfather was a hero in his village.
  \item \emph{Recall:} say the Arabic for a legend, then the Arabic for a myth, then say \emph{baṭal} again
\end{itemize}

\begin{tcolorbox}[breakable,colback=teal!4,colframe=teal!35!black,title={Before you move on}]
What does \ar{بطل} mean? (\textquotedblleft{}A hero; a champion\textquotedblright{}.) Say it once more.
\end{tcolorbox}
\section[\texorpdfstring{\ar{شرير} (shirrīr)}{شرير (shirrīr)}]{\ar{شرير} (shirrīr) — a villain; wicked}

\label{lesson:AR-C303-shirrir}



\begin{quote}
Before the new one: say the Arabic for a myth, then the Arabic for a hero.
\end{quote}

\subsection*{You'll want to know: \ar{شرير}}
\textbf{\ar{شرير}} — \emph{shirrīr} — \textquotedblleft{}a villain; wicked\textquotedblright{}.

\begin{quote}
\ar{كان} \ar{الملك} \ar{في} \ar{الحكاية} \ar{شريرا}.
\end{quote}

\begin{quote}
\emph{kāna al-malik fī al-ḥikāya shirrīran.}
\end{quote}

\begin{quote}
The king in the tale was wicked.
\end{quote}

\begin{cousinweb}
\textbf{\ar{شرير}} (\emph{shirrīr}) is built on the native Arabic root \emph{sh-r-r} (Lane), the root of \emph{sharr}, evil. The \emph{fiʿʿīl} pattern marks a quality held strongly: \textquotedblleft{}very evil\textquotedblright{}.
\end{cousinweb}

\begin{grammarlens}[title={What you've built}]
One more word for telling what happened.
\end{grammarlens}

\subsection*{Your turn}
\begin{itemize}
\raggedright
  \item \emph{Say it:} \emph{shirrīr}
  \item \emph{Say it:} \emph{shirrīr}, once more
  \item \emph{Say it:} \emph{kāna al-malik fī al-ḥikāya shirrīran.} — The king in the tale was wicked.
  \item \emph{Recall:} say the Arabic for a myth, then the Arabic for a hero, then say \emph{shirrīr} again
\end{itemize}

\begin{tcolorbox}[breakable,colback=teal!4,colframe=teal!35!black,title={Before you move on}]
What does \ar{شرير} mean? (\textquotedblleft{}A villain; wicked\textquotedblright{}.) Say it once more.
\end{tcolorbox}
\section[\texorpdfstring{\ar{تخيل} (takhayyala)}{تخيل (takhayyala)}]{\ar{تخيل} (takhayyala) — to imagine}

\label{lesson:AR-C303-takhayyala}



\begin{quote}
Before the new one: say the Arabic for a hero, then the Arabic for a villain.
\end{quote}

\subsection*{You'll want to know: \ar{تخيل}}
\textbf{\ar{تخيل}} — \emph{takhayyala} — \textquotedblleft{}to imagine\textquotedblright{}.

\begin{quote}
\ar{تخيلت} \ar{البنت} \ar{أنها} \ar{تطير} \ar{فوق} \ar{المدينة}.
\end{quote}

\begin{quote}
\emph{takhayyalat al-bint annahā taṭīru fawqa al-madīna.}
\end{quote}

\begin{quote}
The girl imagined that she was flying over the city.
\end{quote}

Present: \ar{يتخيل}, \emph{yatakhayyal}.

It often takes \ar{أن}: \ar{تخيلت} \ar{أنها} …, she imagined that she ….

\begin{cousinweb}
\textbf{\ar{تخيل}} (\emph{takhayyala}) is built on the native Arabic root \emph{kh-y-l} (Lane), \textquotedblleft{}to think, to fancy\textquotedblright{}, the root of \emph{khayāl}, imagination. It is Form V, \textquotedblleft{}to picture to oneself\textquotedblright{}.
\end{cousinweb}

\begin{grammarlens}[title={What you've built}]
One more word for telling what happened.
\end{grammarlens}

\subsection*{Your turn}
\begin{itemize}
\raggedright
  \item \emph{Say it:} \emph{takhayyala}
  \item \emph{Say it:} \emph{takhayyala}, once more
  \item \emph{Say it:} \emph{takhayyalat al-bint annahā taṭīru fawqa al-madīna.} — The girl imagined that she was flying over the city.
  \item \emph{Recall:} say the Arabic for a hero, then the Arabic for a villain, then say \emph{takhayyala} again
\end{itemize}

\begin{tcolorbox}[breakable,colback=teal!4,colframe=teal!35!black,title={Before you move on}]
What does \ar{تخيل} mean? (\textquotedblleft{}To imagine\textquotedblright{}.) Say it once more.
\end{tcolorbox}
